\chapter{Real Analysis: Measuring Volumes on Submanifolds}

\begin{metanote}[Lecture Commencement]
The professor walks to the center of the blackboard, holding a small toy gavel. The blackboard is clean, and the camera is fixed on the main panel.
\end{metanote}

\begin{didacticinsight}[The Gavel and the ``Extra Circus'']
The professor opens the lecture holding a toy gavel, explicitly preparing the students for an ``extra circus''. This playful, theatrical prop serves a distinct pedagogical purpose: acknowledging the escalating difficulty of the material (``The Final Boss of Analysis II'') while keeping the classroom atmosphere grounded and engaged. He beautifully recalls his earlier, informal ``potato'' analogies to show how far the class's rigor has progressed.
\end{didacticinsight}

% ==========================================
% SECTION 1: THE FINAL BOSS (DIVERGENCE THEOREM)
% ==========================================

\section{Recall: The Final Boss of Analysis II}

\begin{spokenclean}[00:00:00 - 00:00:58]
So I brought back the gavel because today I expect extra circus, okay? 

So, remember that I told you on the first day that the final goal, or the final boss of Analysis II, or one of the most representative theorems, is this Divergence Theorem. And in the first days, I needed to state it very informally using ``potatoes'' and stuff because you did not know mathematically any of the elements. But now we saw the statement again some weeks ago, and we already knew some parts. And now we know almost all of it. Right? 
\end{spokenclean}

\begin{nicebox}[Board Action: Divergence Setup]
The professor places the gavel on the desk and begins writing the formal elements of the Divergence Theorem on the left side of the blackboard, underlining the conditions for the domain and the vector field.
\end{nicebox}

\begin{spokenclean}[00:00:58 - 00:01:55]
I will remind you of the statement. Now we know many of the elements: $\Omega$ in $\mathbb{R}^n$ is an open, bounded set. The boundary $\partial \Omega$ will be a $C^1$ submanifold. This we know what it means. Now, $F$ will be a $C^1$ function up to the boundary, taking values in $\mathbb{R}^n$. This is a vector field. 

And the theorem says that the integral over $\Omega$ of this combination of partial derivatives—the derivative $i$ of the component $i$ of the vector field—and you integrate with respect to the Jordan measure (i.e., the Riemann integral, $dx$)... this is equal to this integral over the boundary of the domain with respect to the surface measure of the boundary.
\end{spokenclean}

\begin{spokenclean}[00:01:55 - 00:02:40]
Now we know, and this is $F$ multiplied by the normal vector, which we already saw what is a tangent vector and a normal vector to the boundary. So more or less everything you know what it is now, except for one thing: this integral over a surface, of the differential of volume in the surface (i.e., $dS$). We know how to integrate on the full space, but not over a curved extra surface.
\end{spokenclean}

\begin{orangeformula}[Divergence Theorem (The Final Boss)]
\begin{theorem}[Divergence Theorem]
Let $\Omega \subset \mathbb{R}^n$ be an open, bounded set, and $\partial \Omega$ a $C^1$ submanifold. Let $F \in C^1(\overline{\Omega}, \mathbb{R}^n)$ be a vector field. The theorem states:
\begin{equation} \label{eq:divergence}
\int_\Omega \sum_{i=1}^n \partial_i F_i \, dx = \int_{\partial \Omega} F \cdot \underbrace{\nu}_{\text{Unit vector normal to } \partial \Omega} \, \underbrace{dS}_{\text{Surface measure (Unknown)}}
\end{equation}
\end{theorem}
\end{orangeformula}

\begin{redundantexplanation}[The Unknown Surface Measure]
The entire motivation for this lecture hinges on $dS$. While the left side of the equation uses the standard $n$-dimensional volume measure $dx$ that the students have mastered via the Riemann integral and Fubini's Theorem, the right side requires integrating over an $(n-1)$-dimensional curved manifold embedded in $\mathbb{R}^n$. We do not yet have a mathematical definition for how to rigorously quantify this ``curved volume''.
\end{redundantexplanation}

% ==========================================
% SECTION 2: QUEST LOG - INTEGRATION OVER SUBMANIFOLDS
% ==========================================

\section{Defining Integrals over d-dimensional Submanifolds}

\begin{spokenclean}[00:02:40 - 00:03:35]
And this is where we are going today. We are going to learn, to define what is the integral over the surface. Okay, so next step in our quest: define the integral over $d$-dimensional submanifolds. 

And to define an integral, as it happened in the case of $\mathbb{R}^n$, first we need to define how to measure the volume of sets over a submanifold. So today, towards this goal, what we will do is measure volumes over submanifolds of dimension $d$. 
\end{spokenclean}

\begin{nicebox}[Board Action: Quest Log]
The professor draws a distinct box on the middle panel to represent the current ``Quest Log'', a recurring motif in his lectures to keep the overarching goal visible.
\end{nicebox}

\begin{spokenclean}[00:03:35 - 00:04:20]
Towards this goal, we start by the easiest case. What is the easiest possible submanifold? It is the submanifold of dimension one. Right? So let us start by $d=1$. A submanifold of dimension one, we call it a curve. And we will see how to measure length. So we need to introduce length.
\end{spokenclean}

\begin{mathstroke}[Quest Log: Volumes and Curves]
\textbf{Next Step:} Define $\int$ over $d$-dim submanifolds.

\textbf{Today's Objective:} Measure volumes over submanifolds of $\dim d$.

We build intuition starting with the simplest case:
\begin{itemize}
    \item $\dim d = 1 \implies$ \textbf{Curve} $\implies$ ``Volume'' is defined as \textbf{Length}.
\end{itemize}
\end{mathstroke}

% ==========================================
% SECTION 3: PARAMETRIZED CURVES AND ARC LENGTH
% ==========================================

\section{Formalizing Curves in \texorpdfstring{$\mathbb{R}^n$}{Rn}}    

\begin{spokenclean}[00:04:20 - 00:05:18]
So, how do we rigorously characterize a one-dimensional submanifold? We use what we call a parametrized curve. We start with a simple interval, let us call it $[a, b]$, on the real line, and we map it into our higher-dimensional space $\mathbb{R}^n$ using a function that we will denote as $\gamma$. 

Now, we cannot just use any arbitrary function. We want our curve to be smooth, without any jagged edges or sudden jumps. Therefore, we require $\gamma$ to be continuously differentiable (i.e., $\gamma \in C^1([a, b], \mathbb{R}^n)$).
\end{spokenclean}

\begin{nicebox}[Board Action: Injectivity and Regularity]
He underlines the $C^1$ condition heavily. Then he draws a curve intersecting itself to visually demonstrate why injectivity might be considered, before shifting focus entirely to the non-zero derivative condition to prevent cusps.
\end{nicebox}

\begin{spokenclean}[00:05:18 - 00:06:10]
Furthermore, to avoid pathological cases where the curve stops or creates a sharp cusp, we must enforce that the derivative vector—the velocity—is never the zero vector at any point $t$ (i.e., $\gamma'(t) \neq 0$ for all $t$). If these conditions hold, we say the curve is regular.
\end{spokenclean}

\begin{mathstroke}[Definition of a Regular Parametrized Curve]
We define a \textbf{parametrized curve} as a map:
\[
\gamma: [a, b] \to \mathbb{R}^n
\]
We say that $\gamma$ is a \textbf{regular curve} if it satisfies two conditions:
\begin{enumerate}
    \item $\gamma \in C^1([a, b], \mathbb{R}^n)$ 
    \item $\gamma'(t) \neq 0$ for all $t \in [a, b]$ 
\end{enumerate}
\end{mathstroke}

\begin{spokenclean}[00:06:10 - 00:07:05]
Once we have this regular curve, the immediate next question is: how do we compute its length? Think about physics for a second. If $\gamma(t)$ represents your position in space at time $t$, then the derivative $\gamma'(t)$ is your velocity vector. If you want to know how far you have traveled, you do not just look at the velocity vector; you look at your speed, which is the magnitude of that velocity vector (i.e., $\|\gamma'(t)\|$).
\end{spokenclean}

\begin{spokenclean}[00:07:05 - 00:07:15]
If you integrate your speed over the entire duration of your journey, from time $a$ to time $b$, you obtain the total distance traveled. This total distance is exactly what we define as the length of the curve.
\end{spokenclean}

\begin{orangeformula}[Length of a Parametrized Curve]
\begin{definition}[Arc Length]
Let $\gamma: [a, b] \to \mathbb{R}^n$ be a regular parametrized curve. The \textbf{length} of $\gamma$, denoted as $L(\gamma)$, is defined by the integral of the Euclidean norm of its derivative:
\begin{equation} \label{eq:arc_length}
L(\gamma) = \int_a^b \|\gamma'(t)\| \, dt
\end{equation}
\end{definition}
\end{orangeformula}

% ==========================================
% SECTION 4: THE AXIOMS OF LENGTH
% ==========================================

\section{Geometric Independence and the Axioms of Length}

\begin{spokenclean}[00:07:15 - 00:08:30]
Now, here is the conceptual leap we must make. The length $L(\gamma)$ that we just defined depends explicitly on the parametrization $\gamma$. It depends on the specific function and the specific time interval. But length, fundamentally, is a geometric property of the set of points in space, not a property of how fast you drive along them. 
\end{spokenclean}

\begin{spokenclean}[00:08:30 - 00:09:30]
If two different people drive along the exact same path, but one drives twice as fast, the geometric length of the path they covered is exactly the same, even though their parametrizations are completely different. Therefore, if our formula is to make any geometric sense, we need to prove that it is invariant under reparametrization.
\end{spokenclean}

\begin{mathstroke}[Reparametrization Setup]
Let $\phi: [c, d] \to [a, b]$ be a bijective, $C^1$ function with $\phi'(s) \neq 0$ for all $s \in [c, d]$.
We define a reparametrization of $\gamma$ as:
\[
\tilde{\gamma}(s) = \gamma(\phi(s))
\]

\textbf{Our analytical goal:} We must show that $L(\tilde{\gamma}) = L(\gamma)$.
\end{mathstroke}

\begin{spokenclean}[00:09:30 - 00:10:30]
Before we prove that our integral formula actually survives this reparametrization mathematically, we have to take a step back. We are constructing a measure of volume—specifically, a 1-dimensional volume, which is length. But what does it even mean to be a valid measure of length?
\end{spokenclean}

\begin{spokenclean}[00:10:30 - 00:11:30]
We cannot just write down an integral, compute it, and blindly accept it as ``length''. There are fundamental, axiomatic properties that any concept of length must satisfy in Euclidean space. If our formula fails even one of these, it is not a valid measure. So we must set the ground rules.
\end{spokenclean}

% ==========================================
% SECTION 5: THE FOUR AXIOMS OF LENGTH
% ==========================================

\begin{spokenclean}[00:11:30 - 00:12:30]
If we are going to define a concept as fundamental as length, it cannot just be a random formula we pull from the air. It has to behave in the way our physical universe expects lengths to behave. If it does not, then we have defined something else entirely. 

So, before we even look at the integral of the norm of the velocity, we must ask ourselves: what properties \textit{must} length satisfy? What are the absolute baseline requirements? There will be four. So there will be four fundamental axioms that govern any valid measurement of length.
\end{spokenclean}

\begin{nicebox}[Classroom Interaction: The Catchbox]
The professor walks back to the main desk. He writes the numbers 1 through 4 on the far right blackboard, leaving the definitions blank. He then picks up a bright red throwable ``Catchbox'' microphone and tosses it to a student in the audience.
\end{nicebox}

\begin{spokenclean}[00:12:30 - 00:13:15]
Okay, so what properties must length satisfy? I have written four empty slots on the board. Maybe you can tell me two, and I will complete the other two. Who wants to catch this? Tell me what properties length must satisfy. Who has a guess for the most obvious property a length should have?
\end{spokenclean}

\begin{studentquestion}
It has to be positive? You can't have a negative length. And you just add them together if you have multiple pieces?
\end{studentquestion}

\begin{spokenclean}[00:13:15 - 00:14:15]
Exactly! First, we have positivity, but more rigorously, we need normalization. The length of the straight unit interval (i.e., $[0,1] \times \{0\}$ in $\mathbb{R}^n$) must exactly equal 1. And yes, it must be additive. If we split a curve into two pieces, the total length must be the sum of the lengths of the individual pieces. 
\end{spokenclean}

\begin{spokenclean}[00:14:15 - 00:15:00]
Now I will complete the other two, which are the geometric anchors. We already discussed one: it should be independent of the chosen parametrization. The geometric length does not care how fast you draw the curve. And finally, invariance under Euclidean isometry. If I rotate or translate the curve rigidly in space, its length cannot change.
\end{spokenclean}

\begin{orangeformula}[The Four Axioms of Length]
\begin{definition}[Axioms of Length]
For a definition of length $L(\gamma)$ to be geometrically valid for a curve $\gamma: [a, b] \to \mathbb{R}^n$, it \textit{must} satisfy:
\begin{enumerate}
    \item \textbf{Independence of Parametrization:} $L(\gamma \circ \phi) = L(\gamma)$ for any valid reparametrization $\phi$.
    \item \textbf{Additivity:} For any $c \in (a, b)$, $L(\gamma|_{[a, b]}) = L(\gamma|_{[a, c]}) + L(\gamma|_{[c, b]})$.
    \item \textbf{Euclidean Isometry Invariance:} Rigid rotations and translations do not alter $L(\gamma)$.
    \item \textbf{Normalization:} The standard unit interval $[0,1] \times \{0\}$ has a length of $1$.
\end{enumerate}
\end{definition}
\end{orangeformula}


% ==========================================
% SECTION 6: PROVING REPARAMETRIZATION INVARIANCE
% ==========================================

\section{Proving the Invariance}

\begin{spokenclean}[00:15:00 - 00:16:30]
Let us focus on the first property we established: Independence of Parametrization. We need to rigorously prove that our integral formula actually honors this axiom. 

Suppose we have our reparametrization $\phi$ mapping the interval $[c, d]$ to $[a, b]$. We assume $\phi$ is strictly increasing, so its derivative is strictly positive (i.e., $\phi'(s) > 0$). We want to compute the length of the new, reparametrized curve $\tilde{\gamma}(s) = \gamma(\phi(s))$. 
\end{spokenclean}

\begin{nicebox}[Board Action: The Chain Rule Setup]
The professor walks over to the center panel of the blackboard, pointing at the definition of the regular curve, and begins writing out the composition of the functions to apply the Chain Rule.
\end{nicebox}

\begin{spokenclean}[00:16:30 - 00:17:45]
How do we find its velocity vector? We just apply the Chain Rule. The derivative $\tilde{\gamma}'(s)$ is simply the derivative of the outer function evaluated at the inner function, multiplied by the derivative of the inner function (i.e., $\tilde{\gamma}'(s) = \gamma'(\phi(s)) \cdot \phi'(s)$). 
\end{spokenclean}

\begin{mathstroke}[Applying the Chain Rule]
Given the reparametrized curve $\tilde{\gamma}(s) = \gamma(\phi(s))$ for $s \in [c, d]$.

By the Chain Rule, the velocity vector is:
\[
\tilde{\gamma}'(s) = \gamma'(\phi(s)) \phi'(s)
\]
Taking the Euclidean norm of both sides:
\[
\|\tilde{\gamma}'(s)\| = \|\gamma'(\phi(s)) \phi'(s)\| = \|\gamma'(\phi(s))\| |\phi'(s)|
\]
Since we assume $\phi$ is strictly increasing, $\phi'(s) > 0$, yielding $|\phi'(s)| = \phi'(s)$.
\[
\implies \|\tilde{\gamma}'(s)\| = \|\gamma'(\phi(s))\| \phi'(s)
\]
\end{mathstroke}

\begin{spokenclean}[00:17:45 - 00:19:00]
Now we plug this norm directly into our length formula for $\tilde{\gamma}$. We get the integral from $c$ to $d$ of $\|\gamma'(\phi(s))\| \phi'(s) \, ds$. 

Does this look familiar? It is exactly the setup for the 1-dimensional Change of Variables formula—or as you likely learned it in Analysis I, Integration by Substitution. 
\end{spokenclean}

\begin{spokenclean}[00:19:00 - 00:20:00]
If we substitute $t = \phi(s)$, then our differential $dt$ is exactly $\phi'(s) ds$. The limits of integration change from $c$ and $d$ to $\phi(c)=a$ and $\phi(d)=b$. And as if by magic, the formula collapses back perfectly into the length of the original curve $\gamma$. The geometry is entirely preserved. We have proven the first axiom.
\end{spokenclean}

\begin{mathstroke}[Integration by Substitution]
Computing the length of the reparametrized curve $\tilde{\gamma}$:
\begin{align*}
L(\tilde{\gamma}) &= \int_c^d \|\tilde{\gamma}'(s)\| \, ds \\
&= \int_c^d \|\gamma'(\phi(s))\| \underbrace{\phi'(s) \, ds}_{\text{Substitution: } dt}
\end{align*}

Let $t = \phi(s)$. The differential becomes $dt = \phi'(s) ds$.
The boundaries transform as: $t(c) = \phi(c) = a$ and $t(d) = \phi(d) = b$.

\begin{align*}
L(\tilde{\gamma}) &= \int_a^b \|\gamma'(t)\| \, dt \\
&= L(\gamma)
\end{align*}

\begin{explanationofsteps}
We have successfully demonstrated that the total arc length is invariant under an orientation-preserving reparametrization, confirming the first fundamental axiom of length. The 1D substitution rule elegantly offsets the scaling of the velocity vectors.
\end{explanationofsteps}
\end{mathstroke}